% Shared hierarchical feature path; model-specific layer details are separate.
\definecolor{FigureConnection}{HTML}{4C4D4F}
\begin{tikzpicture}[x=1mm,y=1mm,font=\figurefont\fontsize{8}{10}\selectfont,
 box/.style={draw=FigureBlue,fill=FigureBlueFill,rounded corners=1.8pt,minimum height=12mm,inner sep=1.7mm,align=center},
 style/.style={box,draw=FigureGreen,fill=FigureGreenFill},
 noise/.style={box,draw=FigureYellow,fill=FigureYellowFill,minimum height=9mm},
 wire/.style={-{Latex[length=1.8mm,width=1.2mm]},draw=FigureConnection,line width=.7pt},
 stylewire/.style={wire,draw=FigureGreen},
 noisewire/.style={wire,draw=FigureYellow},
 note/.style={font=\FigureNoteFont\fontsize{7.5}{9.5}\selectfont,text=FigureMuted,align=center}]
\node[box,minimum width=19mm] (z) at (12,72) {随机输入 $\bm z$};
\node[style,minimum width=21mm] (m) at (43,72) {映射网络 $M$};
\node[style,minimum width=19mm] (w) at (75,72) {中间向量 $\bm w$};
\node[box,minimum width=19mm] (c) at (12,36) {可学习常量\\$4\times4$};
\node[box,minimum width=22mm] (b4) at (44,36) {$4\times4$\\合成块};
\node[box,minimum width=25mm] (b8) at (82,36) {$8\times8$\\上采样与合成};
\node[box,minimum width=25mm] (br) at (125,36) {$R\times R$\\上采样与合成};
\node[box,minimum width=18mm] (rgb) at (156,36) {toRGB};
\node[box,minimum width=18mm] (out) at (156,13) {生成图像};
\node[noise,minimum width=22mm] (n4) at (44,10) {独立噪声};
\node[noise,minimum width=22mm] (n8) at (82,10) {独立噪声};
\node[noise,minimum width=22mm] (nr) at (125,10) {独立噪声};
\draw[wire] (z)--(m);
\draw[wire] (m)--(w);
\draw[wire] (c)--(b4);
\draw[wire] (b4)--(b8);
\draw[wire] (b8.east)--(100,36);
\node[inner sep=0,font=\figurefont\fontsize{8}{10}\selectfont] at (103,36) {$\cdots$};
\draw[wire] (106,36)--(br.west);
\draw[wire] (br)--(rgb);
\draw[wire] (rgb)--(out);
\draw[stylewire] (w.south)--(75,58)--(44,58)--(b4.north);
\draw[stylewire] (75,58)--(82,58)--(b8.north);
\draw[stylewire] (82,58)--(125,58)--(br.north);
\node[note] at (115,67) {各层分别得到风格系数};
\draw[noisewire] (n4)--(b4.south);
\draw[noisewire] (n8)--(b8.south);
\draw[noisewire] (nr)--(br.south);
\node[note] at (84,-3) {噪声在块内的各合成层分别注入};
\end{tikzpicture}
